\documentclass[11pt]{article}

\usepackage{acl}

\usepackage{times}
\usepackage{latexsym}
\usepackage[T1]{fontenc}
\usepackage[utf8]{inputenc}
\usepackage{microtype}
\usepackage{inconsolata}
\usepackage{graphicx}
\usepackage{booktabs}
\usepackage{multirow}
\usepackage{amsmath}
\usepackage{enumitem}

\title{Multi-Label Emotion Classification for Bantu Languages:\\
A Comparative Study of TF-IDF and Multilingual Transformer Models}

\author{YM (Malaika) Kamangu, NL (Lebogang) Masenya, DK (David) Musa-Aisien \\
  Group 12\thanks{COS 760 Natural Language Processing --- Department of Computer Science, University of Pretoria --- S1 2026}}

\begin{document}
\maketitle

% ──────────────────────────────────────────────────────────────────────────────
\begin{abstract}
Emotion analysis in African languages remains a significantly under-explored
area despite the continent's remarkable linguistic diversity.
This paper investigates multi-label emotion classification across three
typologically related Bantu languages---isiZulu, isiXhosa, and Swahili---using
the BRIGHTER dataset \citep{abdulmumin2024brighter}.
We compare a TF-IDF with One-vs-Rest Logistic Regression baseline against
two multilingual transformer models, mBERT \citep{devlin2019bert} and
XLM-R \citep{conneau2020unsupervised}, for classifying six emotion categories
(anger, disgust, fear, joy, sadness, surprise).
The baseline achieves a macro F1 of 0.25--0.27 across languages, establishing
a lower bound for the task.
We document substantial class imbalance---particularly for disgust and fear---as
a key challenge, and analyse cross-lingual differences in model behaviour.
Our work contributes a reproducible multi-label emotion pipeline for low-resource
Bantu languages and provides insights into how genealogical relatedness
and data availability influence classifier performance.
\end{abstract}

% ──────────────────────────────────────────────────────────────────────────────
\section{Introduction}

Emotion analysis is a branch of NLP concerned with automatically detecting
affective states expressed in text.
Applications span mental health monitoring, social media analysis, and
opinion mining---areas with growing relevance for the African continent,
where over 2\,000 languages are spoken yet remain largely absent from mainstream
NLP resources \citep{adelani2022masakhaner}.

Research has disproportionately focused on high-resource languages such as
English and Mandarin, leaving African languages underserved in both data
availability and model development.
This resource gap is consequential: without effective tools for low-resource
languages, NLP-based services remain inaccessible to hundreds of millions of
speakers.
Recent community efforts---including AfriSenti \citep{muhammad2023afrisenti}
and AfroXLMR \citep{alabi2022adapting}---demonstrate growing momentum in
African language NLP, yet multi-label emotion classification for Bantu
languages remains largely unaddressed.

The BRIGHTER dataset \citep{abdulmumin2024brighter} addresses this gap by
providing multi-label emotion annotations for 28 languages, including several
Bantu varieties.
In this work, we restrict our experiments to three genealogically related
Bantu languages---isiZulu (\textit{zul}), isiXhosa (\textit{xho}), and
Swahili (\textit{swa})---to investigate whether intra-family linguistic
relatedness influences the efficacy of multilingual pre-trained models.

\paragraph{Research Questions.}
\begin{enumerate}[noitemsep]
  \item How do multilingual transformer models (mBERT, XLM-R) compare to a
        TF-IDF baseline for multi-label emotion classification in Bantu languages?
  \item How does performance vary across isiZulu, isiXhosa, and Swahili, and
        what role does available training data play?
  \item What common prediction errors arise, and what do they reveal about
        class imbalance and model limitations?
\end{enumerate}

\paragraph{Responsible NLP.}
This work directly addresses language underrepresentation.
By building and evaluating pipelines for Bantu languages, we contribute
towards more equitable NLP---recognising that technological inclusion requires
dedicated effort for low-resource communities.
\paragraph{Paper structure.}
Section~2 surveys background and related work.
Section~3 describes the data, models, and evaluation setup.
Section~4 presents experimental results and analysis.
Section~5 reflects on challenges and broader implications.
Section~6 concludes with directions for future work.
% ──────────────────────────────────────────────────────────────────────────────
\section{Background}

\subsection{Multilingual Pre-trained Models}

BERT \citep{devlin2019bert} introduced the transformer encoder architecture
pre-trained with masked language modelling.
Its multilingual variant, mBERT, extends this to 104 languages using a shared
subword vocabulary, enabling zero-shot and few-shot cross-lingual transfer
without language-specific fine-tuning.
XLM-R \citep{conneau2020unsupervised} improves on mBERT through a significantly
larger and more diverse pre-training corpus (2.5 TB of filtered Common Crawl
text across 100 languages), making it particularly effective for low-resource
settings.
Both models represent the current standard for cross-lingual NLP transfer.
However, neither was pre-trained with explicit coverage of Bantu morphological
structure, and their shared subword vocabularies may segment agglutinative
forms in ways that obscure linguistically meaningful boundaries---a gap
directly relevant to isiZulu and isiXhosa.

\subsection{Emotion Analysis in African Languages}

Sentiment and emotion analysis for African languages has grown substantially
with the introduction of benchmarks such as AfriSenti
\citep{muhammad2023afrisenti}, a multilingual Twitter sentiment dataset, and
MasakhaNER \citep{adelani2022masakhaner} for named entity recognition.
\citet{mabokela2024explainable} demonstrated that pre-trained language models
substantially outperform classical approaches for low-resource African language
sentiment, while \citet{thangaraj2024crosslingual} found that XLM-R
consistently outperforms mBERT on cross-lingual transfer for African languages.
However, these studies focus on single-label sentiment or sequence-level tasks;
multi-label emotion classification under severe data constraints (fewer than
900 training examples) has received little systematic attention for Bantu
languages, representing the gap this work addresses.

\subsection{Multi-label Emotion Classification}

Unlike single-label sentiment, multi-label emotion classification assigns
a set of emotion labels to each text instance, reflecting the co-occurrence
of emotions in natural language.
Standard approaches include binary relevance (one classifier per label),
label-powerset methods, and, more recently, sigmoid-activated transformer
classifiers with binary cross-entropy loss
\citep{devlin2019bert, conneau2020unsupervised}.
However, most existing multi-label benchmarks are English-centric; systematic
evaluation under the class imbalance and data scarcity characteristic of
low-resource African languages remains limited.

\subsection{The BRIGHTER Dataset}

BRIGHTER \citep{abdulmumin2024brighter} provides multi-label emotion
annotations across 32 languages, with annotations sourced from fluent speakers.
The dataset includes six emotion categories---anger, disgust, fear, joy,
sadness, and surprise---encoded as binary label vectors, making it directly
applicable to multi-label classification.

\paragraph{Gap.}
Prior work has focused on broad cross-continental transfer without exploiting
intra-family linguistic relatedness.
This project addresses that gap by controlling language selection to a single
genealogical family (Bantu) and examining how data availability interacts with
model performance in a low-resource setting.
The EthioEmo dataset was considered for
Amharic coverage but was not included in experiments due to data availability
constraints; we note this as a limitation and direction for future work.

% ──────────────────────────────────────────────────────────────────────────────
\section{Methodology}

\subsection{Data}

We use the BRIGHTER dataset for three Bantu languages.
Table~\ref{tab:data} summarises split sizes and the emotion label distribution
in the training split for each language.
isiZulu and isiXhosa lack an official training partition; following standard
practice in low-resource NLP, we use the development (validation) split as
training data and evaluate on the held-out test set.

\begin{table}[h]
\centering
\small
\begin{tabular}{lrrrr}
\toprule
\textbf{Lang} & \textbf{Train} & \textbf{Dev} & \textbf{Test} &
  \textbf{Train src.} \\
\midrule
isiZulu  & ---        & 875    & 2{,}047 & validation \\
isiXhosa & ---        & 682    & 1{,}594 & validation \\
Swahili  & 3{,}307    & 1{,}102 & 3{,}312 & train \\
\bottomrule
\end{tabular}
\caption{Dataset split sizes. isiZulu and isiXhosa have no official train
split; the validation set is used as training data.}
\label{tab:data}
\end{table}

\noindent
Label distribution is highly skewed.
For isiXhosa, \textit{sadness} accounts for 41.2\% of training labels while
\textit{disgust} appears in only 0.4\% (3 samples)---a 93.7:1 imbalance ratio.
isiZulu shows moderate imbalance (7.2:1), and Swahili's most severe case is
\textit{fear} at 2.8\% of labels (93 samples, ratio 5.8:1).

\paragraph{Exploratory data analysis.}
In low-resource settings, careful exploratory analysis is essential: limited
data make understanding each sample's structure critical before modelling
\citep{abdulmumin2024brighter}.
Character-level inspection reveals that Swahili text is predominantly
ASCII-compatible, while isiZulu and isiXhosa encode click consonants via
digraphs (\textit{hl}, \textit{xh}, \textit{gc}, \textit{nhl}) and phonemic
modifier letters.
Tokenisation analysis shows that mBERT and XLM-R frequently split these
digraphs into sub-word units that do not align with phonemic boundaries---a
tokenisation challenge inherent to applying models without dedicated Bantu
language coverage.
Minor script inconsistencies (e.g.\ inconsistent apostrophe characters for
click consonants in isiXhosa) were also observed, consistent with the
dialectal variation expected in social-media-sourced annotations.
No cross-lingual alignment was applied; we rely on the models' shared
multilingual subword vocabulary as the sole inter-language bridge.

\subsection{Preprocessing}

Text is used without stemming or lemmatisation, preserving morphological
richness important in Bantu languages.
For the TF-IDF baseline, accents are retained (\texttt{strip\_accents=None})
to avoid losing language-specific characters.
Transformer models apply their native subword tokenisation (WordPiece for
mBERT; SentencePiece for XLM-R) with maximum sequence length 128.

\subsection{Models}

\paragraph{TF-IDF Baseline.}
A TF-IDF vectoriser (maximum 5{,}000 features, unigrams and bigrams,
minimum document frequency 2) feeds a One-vs-Rest Logistic Regression
classifier (\texttt{class\_weight=``balanced''}, \texttt{max\_iter=1000}).
The \texttt{balanced} weighting compensates for label imbalance by inversely
weighting class frequencies.

\paragraph{mBERT.}
\texttt{bert-base-multilingual-cased} \citep{devlin2019bert} is fine-tuned
with a multi-label classification head.
The \texttt{problem\_type=``multi\_label\_classification''} configuration
applies binary cross-entropy with logits (BCEWithLogitsLoss) natively.

\paragraph{XLM-R.}
\texttt{xlm-roberta-base} \citep{conneau2020unsupervised} uses an identical
head and loss configuration to mBERT, enabling direct comparison.

\subsection{Training and Evaluation Setup}

All transformer models are trained under identical, controlled conditions
to ensure fair comparison: 3 epochs, learning rate $2 \times 10^{-5}$,
batch size 8, weight decay 0.01, seed 42.
At inference, logits are converted to probabilities via sigmoid and binarised
at threshold $\tau = 0.3$---a value lower than the standard 0.5 to improve
recall on infrequent emotion classes.

\paragraph{Evaluation metrics.}
Macro F1 is the primary metric, as it treats all classes equally and is robust
to class imbalance.
Supplementary metrics include micro F1, macro precision, macro recall,
Hamming loss, and macro Jaccard score.

% ──────────────────────────────────────────────────────────────────────────────
\section{Experiments and Results}

\subsection{Baseline Performance}

Table~\ref{tab:baseline} reports TF-IDF + Logistic Regression results.
Macro F1 scores range from 0.250 (Swahili) to 0.265 (isiXhosa).
Despite isiXhosa and isiZulu using the development set as training data
(limiting training size to 682 and 875 samples respectively), their macro F1
is comparable to Swahili, which has 3{,}307 training samples.
However, Swahili achieves a notably lower micro F1 (0.297), suggesting its
predictions are more uniformly spread but less precise on frequent classes.

\begin{table}[h]
\centering
\small
\begin{tabular}{lrrrr}
\toprule
\textbf{Language} & \textbf{F1-mac} & \textbf{F1-mic} & \textbf{Prec-mac} &
  \textbf{Rec-mac} \\
\midrule
isiZulu  & 0.257 & 0.341 & 0.235 & 0.290 \\
isiXhosa & 0.265 & 0.471 & 0.247 & 0.321 \\
Swahili  & 0.250 & 0.297 & 0.215 & 0.304 \\
\bottomrule
\end{tabular}
\caption{TF-IDF + Logistic Regression baseline results.}
\label{tab:baseline}
\end{table}

\subsection{Transformer Results}

Tables~\ref{tab:mbert} and~\ref{tab:xlmr} report fine-tuned mBERT and
XLM-R performance.

The mBERT results, fine-tuned with per-label positive-class weighting
(\texttt{pos\_weight}), show improved recall on minority classes compared
to the unweighted configuration.
isiXhosa (682 training samples) achieves macro F1 of 0.240, broadly
comparable to its TF-IDF baseline (0.265), while recall improves markedly
from 0.321 to 0.530 --- indicating that class weighting successfully
pushes predictions towards minority labels.
isiZulu (875 training samples) achieves macro F1 of 0.177, below the
TF-IDF baseline (0.257), but recall rises from 0.290 to 0.560,
confirming that \texttt{pos\_weight} addresses the all-zeros collapse
observed without weighting.
We attribute the remaining macro F1 gap to the data constraint and
the severe sparsity of minority emotion classes (disgust at 2.9\%,
fear at 3.1\% of isiZulu training labels).
Swahili, with the largest training set (3{,}307 examples), achieves macro
F1 of 0.193 --- still below its TF-IDF baseline of 0.250, but with
recall of 0.769, reflecting the weighted loss effectively promoting
minority-class predictions at the cost of precision (0.112).

\begin{table}[h]
\centering
\small
\begin{tabular}{lrrrr}
\toprule
\textbf{Language} & \textbf{F1-mac} & \textbf{F1-mic} & \textbf{Prec-mac} &
  \textbf{Rec-mac} \\
\midrule
isiZulu  & 0.177 & 0.259 & 0.124 & 0.560 \\
isiXhosa & 0.240 & 0.439 & 0.159 & 0.530 \\
Swahili  & 0.193 & 0.232 & 0.112 & 0.769 \\
\bottomrule
\end{tabular}
\caption{mBERT fine-tuning results (3 epochs, batch 8, $\tau{=}0.3$).}
\label{tab:mbert}
\end{table}

\begin{table}[h]
\centering
\small
\begin{tabular}{lrrrr}
\toprule
\textbf{Language} & \textbf{F1-mac} & \textbf{F1-mic} & \textbf{Prec-mac} &
  \textbf{Rec-mac} \\
\midrule
isiZulu  & 0.135 & 0.218 & 0.079 & 0.547 \\
isiXhosa & 0.222 & 0.447 & 0.147 & 0.500 \\
Swahili  & 0.223 & 0.249 & 0.136 & 0.765 \\
\bottomrule
\end{tabular}
\caption{XLM-R fine-tuning results (3 epochs, batch 8, $\tau{=}0.3$).}
\label{tab:xlmr}
\end{table}

XLM-R, also fine-tuned with per-label class weighting, follows a broadly
similar pattern to mBERT.
For isiZulu (875 samples), class weighting prevents the all-zeros collapse
observed in unweighted runs, yielding macro F1 of 0.135 and recall of
0.547; it still falls below mBERT (0.177) owing to XLM-R's larger
parameter count (278M vs.\ mBERT's 178M), which demands more data to
fine-tune effectively.
For isiXhosa (682 samples), XLM-R achieves macro F1 of 0.222, slightly
below mBERT's 0.240, with comparable micro F1 of 0.447.
For Swahili, XLM-R (F1 = 0.223) narrows the gap with mBERT (F1 = 0.193)
when trained on 3{,}307 examples, confirming that XLM-R benefits more
from additional training data.
Notably, the TF-IDF baseline (F1 = 0.250--0.265) outperforms both
transformer models across all three languages under these data constraints,
highlighting the challenge of deploying large pretrained transformers in
genuinely low-resource settings without dedicated hyperparameter tuning
per language.

\subsection{Weak Label Analysis}

Table~\ref{tab:weak} shows the five emotion labels identified as weak
(macro F1 $< 0.15$) in the baseline model.
All weak labels correspond to the rarest classes in their respective
language splits.

\begin{table}[h]
\centering
\small
\begin{tabular}{llrrr}
\toprule
\textbf{Lang} & \textbf{Label} & \textbf{F1} & \textbf{Recall} &
  \textbf{Support} \\
\midrule
isiXhosa & disgust & 0.017 & 0.125 & 8 \\
isiXhosa & fear    & 0.044 & 0.065 & 31 \\
isiXhosa & anger   & 0.069 & 0.091 & 66 \\
isiZulu  & disgust & 0.075 & 0.119 & 59 \\
Swahili  & fear    & 0.069 & 0.064 & 94 \\
\bottomrule
\end{tabular}
\caption{Weak labels (F1 $< 0.15$) from the TF-IDF baseline.
Support reflects test-set samples.}
\label{tab:weak}
\end{table}

\noindent
isiXhosa \textit{disgust} is the most extreme case: despite 8 test samples,
the model achieves only F1 = 0.017, with the extremely limited training
exposure (3 training samples) making generalisation essentially impossible.

\subsection{Error Analysis}

We categorised prediction errors into three types based on the test set:

\begin{itemize}[noitemsep]
  \item \textbf{False negatives}: the model predicts no labels despite ground
        truth labels being present. This is the most common failure mode across
        all three languages.
  \item \textbf{Over-prediction}: the model predicts labels that do not appear
        in the ground truth.
  \item \textbf{Partial match}: some labels are predicted correctly but
        additional labels are missed or spuriously added.
\end{itemize}

False negatives dominate across isiZulu, isiXhosa, and Swahili.
A representative example from isiZulu illustrates the challenge:
the text \textit{``bayaveva ngokujola kukathisha nomfundi''} (expressing anger
about a teacher--student relationship) receives no prediction, likely because
TF-IDF features cannot generalise the infrequent vocabulary.
Over-prediction errors tend to co-occur with sadness and disgust, suggesting
the model associates surface-level negative words with multiple emotion classes
simultaneously.

\paragraph{Responsible NLP.}
Performance disparities across languages are notable: isiXhosa achieves a
higher micro F1 (0.471) than isiZulu (0.341), despite having fewer training
samples. This counterintuitive result is likely driven by isiXhosa's highly
skewed distribution (sadness + joy account for 72.3\% of labels), which makes
predictions on the dominant classes easier.
Such disparities highlight that aggregate metrics can obscure systematic
disadvantage for underrepresented emotion categories, and that evaluation
must be disaggregated per label.

\paragraph{Answers to research questions.}
Returning to the research questions posed in Section~1:
(RQ1) the TF-IDF baseline outperforms both fine-tuned transformer models on
macro F1 across all three languages, demonstrating that larger pre-trained
models are not automatically superior under genuine data scarcity;
(RQ2) performance variation is driven primarily by training data volume and
class distribution rather than linguistic distance---Swahili's larger corpus
yields more balanced precision--recall trade-offs, while isiZulu and isiXhosa
suffer disproportionately from the absence of official training splits;
(RQ3) false negatives are the dominant error type across all languages,
arising from class sparsity that prevents meaningful decision boundaries for
rare emotions such as disgust and fear.

% ──────────────────────────────────────────────────────────────────────────────
\section{Reflections and Discussion}

\subsection{What Worked}

The TF-IDF + Logistic Regression baseline provides a consistent,
reproducible lower bound and is computationally accessible on consumer
hardware.
The \texttt{class\_weight=``balanced''} configuration mitigated---though did
not eliminate---the effect of class imbalance.
The modular pipeline architecture allows baseline and transformer experiments
to share the same data loading and evaluation utilities, ensuring comparability.

\subsection{Challenges}

\paragraph{Limited training data.}
isiZulu and isiXhosa lack official training splits in BRIGHTER, forcing use
of the development set as a training surrogate.
With only 682--875 samples, models are severely data-constrained for a
six-class multi-label problem.

\paragraph{Class imbalance.}
isiXhosa \textit{disgust} appears in 0.4\% of training samples (3 instances),
making it statistically impossible for any model to learn a meaningful
decision boundary.
Even Swahili---with the largest training set---has \textit{fear} at only 2.8\%.
Future work should explore data augmentation techniques, including back-translation
or cross-lingual transfer from higher-resource languages.

\paragraph{Computational constraints.}
Transformer fine-tuning was conducted on CPU, limiting the number of training
epochs feasible within project timelines.
Initial mBERT experiments with a single epoch and batch size of 2 resulted
in training collapse (all-zero predictions), likely due to insufficient gradient
signal before convergence.
Increasing training to 3 epochs, batch size 8, and lowering the classification
threshold to 0.3 addressed this failure mode.

\paragraph{EthioEmo.}
The proposal identified the EthioEmo dataset as a potential Amharic language
complement to the Bantu languages in BRIGHTER.
Due to data availability and time constraints, EthioEmo was not incorporated
into the final experiments.
Including Amharic as a non-Bantu African language reference point remains a
clear direction for future work.

\subsection{Responsible AI Principles}

\paragraph{Data openness.}
All datasets used (BRIGHTER) are publicly available research datasets, and
all code and model artefacts for this project are version-controlled.

\paragraph{Fairness across languages.}
The controlled selection of Bantu languages was an intentional choice to study
cross-lingual fairness within a genealogical family.
Results show that performance inequality exists even among related languages,
driven primarily by training data volume and class distribution rather than
linguistic distance alone.

\paragraph{Limitations and risks.}
The models evaluated here are research prototypes not suitable for deployment.
Emotion classification carries inherent risks of misclassification, which
could have adverse consequences in sensitive applications such as mental health
monitoring.
The near-zero performance on rare emotions (disgust, fear) means these classes
are effectively undetected by all models evaluated.

\paragraph{Broader societal impact.}
Building emotion analysis infrastructure for Bantu languages has positive
long-term implications for social media analysis, public health surveillance,
and community voice amplification.
However, any deployment should involve community consultation and ongoing
fairness auditing.

% ──────────────────────────────────────────────────────────────────────────────
\section{Conclusion}

We presented a comparative evaluation of TF-IDF and multilingual transformer
models for multi-label emotion classification in three Bantu languages
(isiZulu, isiXhosa, Swahili) using the BRIGHTER dataset.

The TF-IDF + Logistic Regression baseline achieves macro F1 of 0.250--0.265,
establishing a reproducible lower bound.
Key findings include: (1) severe class imbalance---particularly for
\textit{disgust} and \textit{fear}---substantially harms recall on minority
classes; (2) false negatives are the dominant error type, suggesting
models default to abstaining rather than over-predicting;
(3) the TF-IDF baseline outperforms both transformer models on macro F1
across all three languages under the data constraints of this study.

For the transformer experiments, fine-tuned with per-label positive-class
weighting (\texttt{pos\_weight}) to address minority-class collapse, mBERT and
XLM-R exhibit complementary data-sensitivity profiles.
mBERT (178M parameters) achieves macro F1 of 0.177 (isiZulu) and 0.240
(isiXhosa), while XLM-R (278M parameters) scores 0.135 and 0.222 respectively,
confirming that the larger model is more data-hungry in low-resource conditions.
Both models substantially improve recall over the baseline (mBERT recall: 0.530--0.769
vs.\ baseline 0.290--0.321), but at the cost of precision, reflecting the
pos\_weight trade-off between sensitivity and specificity.
For Swahili (3{,}307 training examples), XLM-R (F1 = 0.223) narrows the gap
with mBERT (F1 = 0.193), consistent with XLM-R benefiting more from larger
training sets.

These findings underscore that data volume is the primary bottleneck for
emotion analysis in low-resource Bantu languages, and that larger pretrained
models are not automatically superior when training data is scarce.
Future work should investigate data augmentation, AfroXLMR
\citep{alabi2022adapting}, and cross-lingual transfer from richer
corpora to address the fundamental data scarcity that constrains all
models evaluated in this setting.

% ──────────────────────────────────────────────────────────────────────────────
\bibliography{custom}

\end{document}
